% Rolling-Shutter Vibrometry

\documentclass[10pt,twocolumn,letterpaper]{article}

% Geometry and fonts
\usepackage[
  paperwidth=8.5in, paperheight=11in,
  left=0.75in, right=0.75in,
  top=1in, bottom=1in
]{geometry}
\usepackage{times}
\usepackage{graphicx}
\usepackage{amsmath,amssymb}
\usepackage{booktabs}
\usepackage{caption}
\usepackage{subcaption}
\usepackage[breaklinks=true,bookmarks=false]{hyperref}
\usepackage{microtype}
\usepackage{xcolor}
\usepackage{enumitem}
\usepackage{fancyhdr}
\usepackage{titlesec}
\usepackage{float} 

\setlength{\columnsep}{0.375in}

% Section headings
\titleformat{\section}{\normalfont\large\bfseries}{\thesection}{0.5em}{}
\titleformat{\subsection}{\normalfont\normalsize\bfseries}{\thesubsection}{0.5em}{}
\titlespacing*{\section}{0pt}{8pt plus 2pt minus 2pt}{4pt plus 1pt minus 1pt}
\titlespacing*{\subsection}{0pt}{6pt plus 2pt minus 2pt}{3pt plus 1pt minus 1pt}

% Convenience macros
\newcommand{\etal}{\textit{et al.}}
\DeclareMathOperator*{\argmax}{arg\,max}

% Page style - no page numbers
\pagestyle{empty}

% Title block
\title{\vspace{-1em}%
  \textbf{Rolling-Shutter Vibrometry:}\\
  Recovering Frequencies from Video Without a Microphone%
  \vspace{-0.5em}
}
\author{%
  Emmanuel Kasuti Makau\\
  Computational Photography, Spring 2026\\
  \texttt{Dartmouth College}
}
\date{}

\begin{document}
\maketitle
\thispagestyle{empty}

% Abstract
\begin{abstract}
Most phones and cameras do not photograph the entire frame at once. Instead, they read the image one row of pixels at a time, from top to bottom. This is called a rolling shutter. Because each row is captured a tiny fraction of a second after the row above it, a vibrating object photographed by such a sensor gets smeared into a wavy pattern: each frame quietly records how the object moved during the short time the sensor took to scan it. I use this normally unwanted effect on purpose. By watching how a vibrating object wobbles inside a single video frame, I can recover the frequency of the sound driving it without a microphone and without special equipment. I stick a strip of tape on a speaker cone to give a sharp edge to follow. For every row of the image I find where that edge sits, clean up the measurement with a median filter, combine several frames to sharpen the frequency reading, and read off the strongest peak in the frequency spectrum. After a single one-time calibration clip, I recover a 50\,Hz tone with 0\% error from an iPhone 16 Pro Max filming in 240\,fps slow-motion, and a 440\,Hz tone with 0\% error from a Canon Rebel T5i DSLR at 30\,fps. I measure a signal-to-noise ratio of 44.4\,dB for the 440\,Hz case, and explain two physical limits on how far the method generalizes: the speaker's own mechanical resonance, and how finely a single rolling-shutter frame can resolve frequencies.
\end{abstract}

% 1. Introduction
\section{Introduction}

Sound measurement ordinarily requires a dedicated device, a microphone or accelerometer, that turns pressure waves in the air into an electrical signal. This constraint can be limiting in surveillance, remote sensing, and accessibility applications where deploying additional hardware is impractical. Yet almost every environment already contains cameras. If an ordinary camera could double as a sound sensor, a whole range of problems suddenly becomes solvable with gear that is already in place.

The key observation is that everyday CMOS camera sensors do not capture all rows of a frame at the same instant. Rows are read top-to-bottom at a fixed rate; the time gap between the first and last row of a 1080-line sensor is on the order of $1/f_\text{readout}$ seconds. When a speaker cone or any vibrating surface is imaged, its position shifts between row reads, producing a wobble of a vertical edge across the frame (Figure~\ref{fig:edge}).

This paper makes the following contributions:
\begin{itemize}[leftmargin=1.2em,itemsep=1pt,topsep=2pt]
  \item A self-contained pipeline for rolling-shutter vibrometry, written in Python.
  \item Experiments across four tones (50 - 1200\,Hz) under two filming conditions: slow-motion iPhone footage and ordinary 30\,fps DSLR footage.
  \item An account of where the method breaks down, speaker mechanical resonance and frequency-resolution limits, with ideas for fixing each one.
\end{itemize}

% 2. Related Work
\section{Related Work}

\paragraph{Visual Microphone.}
Davis~\etal~\cite{davis2014} demonstrated that subtle, sub-pixel vibrations of everyday objects (a crisp packet, a plant) can be recovered from high-speed video by analyzing motion in localized regions. Their method relies on Eulerian motion magnification~\cite{wu2012} and requires either a high-speed camera or a carefully designed optical-flow pipeline. My work tackles the opposite end of the spectrum: standard or slow-motion smartphone footage, in which the rolling-shutter effect becomes the primary sensing mechanism, with no motion magnification required.

\paragraph{Rolling-Shutter Imaging.}
The geometric distortions caused by rolling shutters have been studied extensively, mostly with the goal of removing them~\cite{forssen2010,grundmann2012}. More recently, researchers have exploited rolling shutter to work as a coded-aperture trick for estimating depth~\cite{albl2016} and for measuring the speed of fast-moving objects~\cite{ringaby2012}. Using rolling shutter directly as a sound sensor is, to my knowledge, novel in tracking the edge position row by row.

\paragraph{Frequency Estimation from Video.}
Poh~\etal~\cite{poh2010} extract heart rate from facial video by detecting subtle color changes due to blood flow. Analogously, I detect mechanical frequency from spatial pixel displacements. The key difference is that my signal is deterministic (a pure tone) and confined to a controlled region (the tape edge), making FFT-based estimation natural and interpretable.

% 3. Method
\section{Method}
\label{sec:method}

\subsection{Physical Model}

Let $f_s$ denote the sensor row readout rate (rows\,s$^{-1}$) and $f_v$ the vibration frequency of the speaker cone. The cone displacement at row $r$ is

\begin{equation}
  x(r) = A\sin\!\left(2\pi f_v \tfrac{r}{f_s} + \phi\right),
  \label{eq:model}
\end{equation}

where $A$ is the amplitude in pixels and $\phi$ is an arbitrary phase. The tape edge appears at column $c_0 + x(r)$ in row $r$. Taking the Fourier transform of the sequence $\{x(r)\}_{r=0}^{R-1}$ yields a spectral peak at bin $k^* = \lfloor f_v R / f_s \rceil$.

\subsection{Pipeline}

\textbf{Step1 - Frame extraction.}
I extract evenly spaced frames from each video clip using OpenCV, saving as 85\% quality JPEG at $1920\times1080$ (trial1, DSLR) or $1920\times1080$ slow-motion (trial5, iPhone).

\textbf{Step 2 - Edge detection.}
For each frame I crop to columns $[950,\,1150]$, a 200-pixel window centered on the tape edge. The per-row edge position is the column of maximum absolute horizontal gradient:
\begin{equation}
  \hat{c}(r) = \argmax_{c}
    \bigl|\,I(r,c+1) - I(r,c)\,\bigr|.
  \label{eq:edge}
\end{equation}
A median filter (kernel size 15) removes spike outliers caused by the speaker screws and rim (Figure~\ref{fig:edge}).

\textbf{Step 3 - Multi-frame stacking.}
A single frame of $H=1080$ rows yields frequency resolution $\Delta f = f_s / H$. Concatenating $K$ frames after per-frame mean subtraction yields $K\!H$ samples and $\Delta f = f_s / (KH)$. I use $K=10$ for trial5 ($\Delta f \approx 5$\,Hz/bin) and $K=10$ for trial1 ($\Delta f \approx 44$\,Hz/bin).

\textbf{Step 4 - FFT and peak detection.}
I apply the real-valued FFT to the concatenated, de-meaned signal and find
the dominant bin (excluding DC):
\begin{equation}
  \hat{f}_v = \argmax_{k>0}
    \bigl|\mathcal{F}\{x\}[k]\bigr| \cdot \tfrac{f_s}{N},
  \label{eq:fft}
\end{equation}
where $N = KH$ is the total sample count.

\subsection{Calibration}
\label{sec:cal}

The nominal row readout rate $f_s^{(0)} = \mathrm{FPS}\times H$ differs from the true physical rate because camera firmware does not expose this parameter directly. I calibrated once with a clip of known frequency $f_\text{ref}$:

\begin{equation}
  f_s = f_s^{(0)} \cdot \frac{f_\text{ref}}{\hat{f}_\text{raw}},
  \label{eq:cal}
\end{equation}

where $\hat{f}_\text{raw}$ is the FFT peak under the nominal rate. For trial5 (iPhone, 240\,fps), using $f_\text{ref}=50$\,Hz yields $f_s = 54{,}000$\,rows\,s$^{-1}$. For trial1 (Canon, 30\,fps), using $f_\text{ref}=440$\,Hz yields $f_s = 475{,}200$\,rows\,s$^{-1}$ — a $14.7\times$ correction factor explained by the sensor's true line period rather than the declared frame rate.

% 4. Experimental Setup
\section{Experimental Setup}

\textbf{Hardware.}
I drove a passive speaker at exact frequencies using a tone generator (\texttt{onlinetonegenerator.com}). A strip of white electrical tape fixed to the speaker cone gave the camera a sharp, high-contrast vertical edge to track. The camera sat on a tripod aimed directly at the edge; only the speaker cone moved. Camera shake adds noise to the edge signal and raises the noise floor. Figure~\ref{fig:setup} shows the complete physical arrangement and the tape displacement pattern it produces.

\begin{figure}[h]
    \centering
    \includegraphics[width=\linewidth]{fig_tape_movement_panel.png}
    \caption{
        \textbf{Setup.} Tape on the speaker cone (left) forms the tracked edge; at 50\,Hz it oscillates row-by-row within one frame (right).}
    \label{fig:setup}
\end{figure}

\textbf{Audio verification.}
Before filming each tone I confirmed the speaker was actually producing the target frequency. Figure~\ref{fig:audio_verify} shows the audio spectrum recorded by a separate microphone: the dominant peak matches the intended tone, ruling out a mismatch between the tone-generator setting and the acoustic output as a confound in the accuracy results.

\begin{figure}[t]
  \centering
   \includegraphics[width=\linewidth]{fig_audio_verification.png}
   \caption{
    \textbf{Audio verification.} Microphone spectrum confirms the speaker   outputs the intended tone before each filming session.}
   \label{fig:audio_verify}
\end{figure}

\textbf{Trial5 (iPhone slow-motion)}
Camera: iPhone 16 Pro Max, 1920$\times$1080, declared 240\,fps.
Calibrated $f_s = 54{,}000$\,rows\,s$^{-1}$.
Test tones: 50, 60, 70\,Hz.

\textbf{Trial1-4 (Canon 30\,fps)}
Camera: Canon Rebel T5i DSLR, 1920$\times$1080, 30\,fps.
Calibrated $f_s = 475{,}200$\,rows\,s$^{-1}$.
Test tones: 200, 440, 880, 1200\,Hz.

\textbf{Silence baseline.}
A clip recorded with no tone playing provides the noise floor for SNR measurement and confirms that any edge oscillation or wobble is caused by the speaker.

% 5. Results
\section{Results}

\subsection{Baseline Verification}

Figure~\ref{fig:baseline} shows the per-row edge position for the silence clip (top) and the 50\,Hz clip (bottom) from the iPhone slow-motion capture. The silence trace is nearly flat; the 50\,Hz trace oscillates visibly, confirming the displacement is driven by sound. Figure~\ref{fig:edge_zoom} shows the same two signals zoomed to 200 rows so the sinusoidal shape is unambiguous.

\begin{figure}[t]
    \centering
    \includegraphics[width=\linewidth]{fig_baseline_silence_vs_50hz.png}
    \caption{
        \textbf{Silence vs.\ 50\,Hz (iPhone).} No tone: flat edge (top). 50\,Hz: clear sinusoidal oscillation (bottom).}
    \label{fig:baseline}
\end{figure}

\begin{figure}[t]
  \centering
  \includegraphics[width=\linewidth]{fig_edge_zoomed_silence_vs_50hz.png}
  \caption{
    \textbf{Zoomed — silence vs.\ 50\,Hz.} Over 200 rows: silence is flat (top), 50\,Hz completes one full cycle (bottom).}
  \label{fig:edge_zoom}
\end{figure}

\subsection{Edge Signal and Frequency Spectrum}

Figure~\ref{fig:edge} shows the raw and filtered edge position for one frame of the 50\,Hz clip. The sinusoidal modulation is clearly visible. A single 1080-row frame gives coarse frequency resolution ($\Delta f \approx 50$\,Hz before calibration). Figure~\ref{fig:fft_single} shows the FFT of just one frame: the correct bin wins, but neighboring bins carry significant energy. Stacking ten frames sharpens the peak enough to eliminate that ambiguity.

\begin{figure}[t]
  \centering
  \includegraphics[width=\linewidth]{fig_edge_signal_raw_vs_smooth.png}
  \caption{
    \textbf{Edge signal — 50\,Hz.} Raw (gray) oscillates with vibration; median filter (red) removes screws/rim outliers.}
  \label{fig:edge}
\end{figure}

\begin{figure}[t]
  \centering
  \includegraphics[width=\linewidth]{fig_fft_single_frame_50hz.png}
  \caption{
    \textbf{Single-frame FFT — 50\,Hz.} Correct peak found, but coarse resolution leaks energy into neighboring bins; stacking 10 frames fixes this.}
  \label{fig:fft_single}
\end{figure}

\subsection{Calibration Verification}
\label{sec:cal_verify}

The calibration step (Section~\ref{sec:cal}) corrects the nominal row readout rate using a clip of known frequency. Figure~\ref{fig:fft_cal} shows the post-calibration FFT for the iPhone at 50\,Hz: the peak lands exactly on the correct bin. Figure~\ref{fig:val440} shows the same check for the Canon at 440\,Hz. In both cases the calibrated estimate matches ground truth to within one FFT bin, validating Equation~\ref{eq:cal}.

\begin{figure}[t]
  \centering
  \includegraphics[width=\linewidth]{fig_fft_calibrated_50hz.png}
  \caption{
    \textbf{Calibrated FFT — 50\,Hz (iPhone).} Peak lands exactly on 50\,Hz (red dashed = ground truth) after $f_s$ correction.}
  \label{fig:fft_cal}
\end{figure}

\begin{figure}[t]
  \centering
  \includegraphics[width=\linewidth]{fig_validation_440hz.png}
  \caption{
    \textbf{Validation — 440\,Hz (Canon).} Clean narrow peak at 440\,Hz confirms the $14.7\times$ $f_s$ correction is accurate.}
  \label{fig:val440}
\end{figure}

\subsection{Frequency Estimation Accuracy}

Table~\ref{tab:results} reports the ground-truth tone, the video-only FFT estimate, and the relative error for every test clip. Figure~\ref{fig:fft_iphone} shows the three FFT spectra from the iPhone slow-motion capture; Figure~\ref{fig:trial1_acc} shows the accuracy summary for the Canon DSLR.

\begin{table}[h]
    \centering
    \caption{Frequency recovery results. GT = ground-truth tone frequency.}
    \label{tab:results}
    \small
    \setlength{\tabcolsep}{4pt}
    \begin{tabular}{lrrr}
        \toprule
            Setup & GT (Hz) & Est. (Hz) & Error (\%) \\
        \midrule
            iPhone 240\,fps & 50  &  50.0 & \textbf{0.0} \\
            iPhone 240\,fps & 60  &  50.0 & 16.7 \\
            iPhone 240\,fps & 70  &  50.0 & 28.6 \\
        \midrule
            Canon 30\,fps   & 200 & 440.0 & 120.0 \\
            Canon 30\,fps   & 440 & 440.0 & \textbf{0.0} \\
            Canon 30\,fps   & 880 & 440.0 &  50.0 \\
            Canon 30\,fps   &1200 & 440.0 &  63.3 \\
        \bottomrule
    \end{tabular}
\end{table}

\begin{figure}[t]
  \centering
  \includegraphics[width=\linewidth]{fig_fft_trial5_all.png}
  \caption{
    \textbf{FFT — iPhone (240\,fps).} 50\,Hz recovered exactly; 60 and 70\,Hz both lock onto the speaker's 50\,Hz resonance.}
  \label{fig:fft_iphone}
\end{figure}

\begin{figure}[t]
  \centering
  \includegraphics[width=\linewidth]{fig_trial1_accuracy.png}
  \caption{
    \textbf{Accuracy — Canon (30\,fps).} True vs.\ estimated frequency (left) and per-tone error (right); 440\,Hz is the only successful recovery.}
  \label{fig:trial1_acc}
\end{figure}

The method achieves \textbf{0\% error} on the two calibration anchors (50\,Hz iPhone, 440\,Hz Canon). These are genuine cross-validation results: calibration uses a single-frame analysis of the anchor clip while accuracy testing uses a separate 10-frame stack, so the correct recovery is not a trivial artifact of the calibration data.

\subsection{Signal-to-Noise Ratio}

Figure~\ref{fig:snr} compares the FFT spectrum of the 440\,Hz clip against the silence baseline. The signal peak rises \textbf{44.4\,dB} above the noise floor, providing a comfortable margin for reliable peak detection with no thresholding or post-processing required.

\begin{figure}[h]
  \centering
  \includegraphics[width=\linewidth]{fig_snr.png}
  \caption{
    \textbf{SNR — Canon (30\,fps).} 440\,Hz signal (orange) rises 44.4\,dB above the silence floor (blue).}
  \label{fig:snr}
\end{figure}

\clearpage
% 6. Limitations and Future Work
\section{Limitations and Future Work}
\label{sec:limits}

\paragraph{Speaker mechanical resonance.}
A speaker cone vibrates most strongly at its natural resonant frequency regardless of the electrical drive. In the iPhone slow-motion capture the cone resonates at $\approx$50\,Hz; when driven at 60 or 70\,Hz the rolling-shutter signal is dominated by that resonance, so the FFT peaks at 50\,Hz regardless of input. This is a \emph{transducer} constraint. Two mitigations are (i)~attaching a rigid retroreflector directly to the cone to decouple mechanical resonance from optical displacement, and (ii)~using a tweeter whose resonant frequency lies above the measurement range.

\paragraph{Calibration-anchor aliasing.}
For the Canon DSLR, calibrating with 440\,Hz scales $f_s$ so the anchor falls at exactly one spatial cycle per frame. Weaker harmonic and sub-harmonic components (880 and 200\,Hz) are masked by residual energy at the calibration bin. Calibrating with a non-resonant anchor and applying sub-pixel interpolation would reduce this masking.

% 7. Conclusion
\section{Conclusion}
I have shown that the rolling-shutter effect present in every commodity camera can be repurposed as a sound sensor without any additional hardware. The pipeline recovers tones from video with zero error at the calibration anchors (50\,Hz from an iPhone 16 Pro Max in slow-motion; 440\,Hz from a Canon Rebel T5i at 30\,fps) and a signal-to-noise ratio of 44.4\,dB. It fails for reasons I understand physically and can fix with hardware changes, which makes it a solid starting point for further work on silent, video-only sound sensing.

% References
{\small
\begin{thebibliography}{9}

\bibitem{davis2014}
A.~Davis, M.~Rubinstein, N.~Wadhwa, G.~Mysore, F.~Durand, and W.~T.~Freeman,
``The Visual Microphone: Passive Recovery of Sound from Video,''
\textit{ACM Trans.\ Graph.\ (SIGGRAPH)}, 33(4):79, 2014.

\bibitem{wu2012}
H.-Y.~Wu, M.~Rubinstein, E.~Shih, J.~Guttag, F.~Durand, and W.~T.~Freeman,
``Eulerian Video Magnification for Revealing Subtle Changes in the World,''
\textit{ACM Trans.\ Graph.\ (SIGGRAPH)}, 31(4):65, 2012.

\bibitem{forssen2010}
P.-E.~Forss\'{e}n and D.~Ringaby,
``Rectifying rolling shutter video from hand-held devices,''
in \textit{CVPR}, 2010.

\bibitem{grundmann2012}
M.~Grundmann, V.~Kwatra, D.~Castro, and I.~Essa,
``Calibration-free rolling shutter removal,''
in \textit{ICCP}, 2012.

\bibitem{albl2016}
C.~Albl, Z.~Kukelova, and T.~Pajdla,
``R6P --- Rolling shutter absolute camera pose,''
in \textit{CVPR}, 2016.

\bibitem{ringaby2012}
E.~Ringaby and P.-E.~Forss\'{e}n,
``Efficient video rectification and stabilisation for cell-phones,''
\textit{IJCV}, 96(3):335--352, 2012.

\bibitem{poh2010}
M.-Z.~Poh, D.~J.~McDuff, and R.~W.~Picard,
``Non-contact, automated cardiac pulse measurements using video imaging and
blind source separation,''
\textit{Opt.\ Express}, 18(10):10762--10774, 2010.

\end{thebibliography}
}

\end{document}
